%% Part 2.2 normal-form fragment: v1.1 Part III, lines 5135--5566.
%% Preserves R-007/R-008.  Apparatus: A-033, A-034, A-118; depth D-P3-012.

\chapter{First-Order Normal Forms and Clausification}
\label{sec:p02-first-order-normal-forms}
\chaptermark{First-Order Normal Forms}

\section{Negation normal form and quantifier laws}
\label{sec:p02-fol-nnf}

\begin{hbdefinition}{def:p02-fol-literal-nnf}{First-order literal and negation normal form}
A \emph{first-order literal}\index{literal!first-order} is an atomic formula
or the negation of one.  A formula is in \emph{negation normal form}
(NNF)\index{negation normal form!first-order} iff it is generated from
first-order literals, \(\bottom\), and \(\neg\bottom\) using only
\(\land,\lor,\forall,\exists\).  Thus every other negation has been pushed
through both propositional connectives and quantifiers.  This is a syntactic
class of formulas, not an assertion of equivalence.
\end{hbdefinition}

\begin{hbmethod}{def:p02-fol-nnf-transform}{Paired first-order NNF recursion}
Define metalinguistic maps
\index{negation normal form!transformation}
\(\operatorname{nnf},\operatorname{nnf}^{\neg}:
\mathsf{Fm}(\mathcal L)\to\mathsf{Fm}(\mathcal L)\) simultaneously.
For an atom \(A\), put
\[
 \operatorname{nnf}(A)=A,
 \quad \operatorname{nnf}^{\neg}(A)=\neg A,
 \quad \operatorname{nnf}(\bottom)=\bottom,
 \quad \operatorname{nnf}^{\neg}(\bottom)=\neg\bottom.
\]
The remaining recursive clauses are
\[
\begin{aligned}
 \operatorname{nnf}(\neg\varphi)
   &=\operatorname{nnf}^{\neg}(\varphi),&
 \operatorname{nnf}^{\neg}(\neg\varphi)
   &=\operatorname{nnf}(\varphi),\\
 \operatorname{nnf}(\varphi\land\psi)
   &=\operatorname{nnf}(\varphi)\land\operatorname{nnf}(\psi),&
 \operatorname{nnf}^{\neg}(\varphi\land\psi)
   &=\operatorname{nnf}^{\neg}(\varphi)\lor
     \operatorname{nnf}^{\neg}(\psi),\\
 \operatorname{nnf}(\varphi\lor\psi)
   &=\operatorname{nnf}(\varphi)\lor\operatorname{nnf}(\psi),&
 \operatorname{nnf}^{\neg}(\varphi\lor\psi)
   &=\operatorname{nnf}^{\neg}(\varphi)\land
     \operatorname{nnf}^{\neg}(\psi),\\
 \operatorname{nnf}(\varphi\to\psi)
   &=\operatorname{nnf}^{\neg}(\varphi)\lor\operatorname{nnf}(\psi),&
 \operatorname{nnf}^{\neg}(\varphi\to\psi)
   &=\operatorname{nnf}(\varphi)\land
     \operatorname{nnf}^{\neg}(\psi),\\
 \operatorname{nnf}(\forall x\,\varphi)
   &=\forall x\,\operatorname{nnf}(\varphi),&
 \operatorname{nnf}^{\neg}(\forall x\,\varphi)
   &=\exists x\,\operatorname{nnf}^{\neg}(\varphi),\\
 \operatorname{nnf}(\exists x\,\varphi)
   &=\exists x\,\operatorname{nnf}(\varphi),&
 \operatorname{nnf}^{\neg}(\exists x\,\varphi)
   &=\forall x\,\operatorname{nnf}^{\neg}(\varphi).
\end{aligned}
\]
Every recursive call is on a formula of smaller complexity.  The companion
map is necessary because it records the NNF of a negation without introducing
a third syntactic operation; both maps are used in
Theorem~\ref{thm:p02-fol-nnf} and the pipeline below.
\end{hbmethod}

\begin{hblemma}{thm:p02-fol-nnf}{proved=annex:p2.fol.nf.1}{First-order NNF theorem}
For every \(\varphi\in\mathsf{Fm}(\mathcal L)\), both
\(\operatorname{nnf}(\varphi)\) and
\(\operatorname{nnf}^{\neg}(\varphi)\) are in NNF,
\[
 \operatorname{nnf}(\varphi)\Leftrightarrow\varphi,
 \qquad
 \operatorname{nnf}^{\neg}(\varphi)\Leftrightarrow\neg\varphi,
\]
and each output has exactly the same free-variable set and all-variable set
as \(\varphi\).
\end{hblemma}

\begin{hbtheorem}{thm:p02-fol-quantifier-equivalences}{cited=Enderton2001}{Quantifier equivalence laws}
Every row of Table~\ref{tab:p02-quantifier-equivalences} is valid in the
nonempty-domain first-order semantics.  A row carrying
\(x\notin\FV(\chi)\) may be used only under that displayed side condition.
\end{hbtheorem}

\begingroup\small\setlength{\tabcolsep}{4pt}
\begin{longtable}{@{}p{.32in}p{2.55in}p{1.88in}@{}}
\caption{Quantifier equivalence laws with side conditions (apparatus A-033).}
\label{tab:p02-quantifier-equivalences}\\
\toprule
Row & Equivalence & Side condition \\ \midrule
\endfirsthead
\toprule Row & Equivalence & Side condition \\ \midrule
\endhead
Q1 & \(\neg\forall x\,\varphi\Leftrightarrow
       \exists x\,\neg\varphi\) & none \\
Q2 & \(\neg\exists x\,\varphi\Leftrightarrow
       \forall x\,\neg\varphi\) & none \\
Q3 & \((Qx\,\varphi)\land\chi\Leftrightarrow
       Qx(\varphi\land\chi)\), \(Q\in\{\forall,\exists\}\) &
       \(x\notin\FV(\chi)\) \\
Q4 & \((Qx\,\varphi)\lor\chi\Leftrightarrow
       Qx(\varphi\lor\chi)\), \(Q\in\{\forall,\exists\}\) &
       \(x\notin\FV(\chi)\) \\
Q5 & \((\forall x\,\varphi)\to\chi\Leftrightarrow
       \exists x(\varphi\to\chi)\) & \(x\notin\FV(\chi)\) \\
Q6 & \((\exists x\,\varphi)\to\chi\Leftrightarrow
       \forall x(\varphi\to\chi)\) & \(x\notin\FV(\chi)\) \\
Q7 & \(\chi\to(\forall x\,\varphi)\Leftrightarrow
       \forall x(\chi\to\varphi)\) & \(x\notin\FV(\chi)\) \\
Q8 & \(\chi\to(\exists x\,\varphi)\Leftrightarrow
       \exists x(\chi\to\varphi)\) & \(x\notin\FV(\chi)\) \\
Q9 & \(\forall x(\varphi\land\psi)\Leftrightarrow
       (\forall x\varphi)\land(\forall x\psi)\) & none \\
Q10 & \(\exists x(\varphi\lor\psi)\Leftrightarrow
         (\exists x\varphi)\lor(\exists x\psi)\) & none \\
Q11 & \(\forall x\,\varphi\Leftrightarrow\varphi\) &
        \(x\notin\FV(\varphi)\) \\
Q12 & \(\exists x\,\varphi\Leftrightarrow\varphi\) &
        \(x\notin\FV(\varphi)\) \\
\bottomrule
\end{longtable}
\endgroup

The converses of the familiar one-way distribution laws omitted from the
table are generally invalid; for example, two existential conjuncts need not
share a witness.  The side conditions in Q3--Q8 are part of the law, not
typographical annotations.

\section{Rectification and prenex form}
\label{sec:p02-fol-prenex}

\begin{hbdefinition}{def:p02-prenex-rectified}{Prenex and rectified formulas}
A formula is \emph{prenex}\index{prenex form} iff it has the form
\[
 Q_1x_1\cdots Q_nx_n\,\chi,
\]
where each \(Q_i\in\{\forall,\exists\}\) and \(\chi\) is
quantifier-free.  It is \emph{rectified}\index{rectification} iff no
variable has both a free and a bound occurrence and distinct quantifier
occurrences bind distinct variables.
\end{hbdefinition}

\begin{hblemma}{thm:p02-fol-rectification}{proved=annex:p2.fol.nf.2}{Rectified variant}
For every formula \(\varphi\) and finite
\(V\supseteq\operatorname{Var}(\varphi)\), define the metalinguistic map
\(\operatorname{rect}\) recursively by fixing atoms and \(\bottom\), commuting
with negation, and putting
\[
\begin{aligned}
 \operatorname{rect}(\psi\mathbin\circ\chi,V)
   &=\operatorname{rect}(\psi,V)\mathbin\circ
     \operatorname{rect}\bigl(\chi,
       V\cup\operatorname{Var}(\operatorname{rect}(\psi,V))\bigr),\\
 \operatorname{rect}(Qx\,\psi,V)
   &=Qz\,\operatorname{rect}(\psi[z/x],V\cup\{z\}),
\end{aligned}
\]
where \(\circ\in\{\land,\lor,\to\}\), \(Q\in\{\forall,\exists\}\), and
\(z\) is the least variable outside \(V\).  This map selects a definite
rectified representative; it is not the alpha-equivalence relation itself and
is needed as the input discipline for \(\operatorname{pren}\).  Then
\begin{enumerate}
  \item \(\operatorname{rect}(\varphi,V)=_{\alpha}\varphi\), hence
        \(\operatorname{rect}(\varphi,V)\Leftrightarrow\varphi\);
  \item it is rectified, its pairwise-distinct binders lie outside \(V\), and
        \[
         \FV(\operatorname{rect}(\varphi,V))=\FV(\varphi);
        \]
  \item it preserves the constructor tree, so it preserves NNF.
\end{enumerate}
Alpha-equivalent formulas have equal free-variable sets.
\end{hblemma}

\begin{hblemma}{thm:p02-fol-prenex-extraction}{proved=annex:p2.fol.nf.3}{Prenex extraction}
For rectified NNF formulas, define the metalinguistic map
\(\operatorname{pren}\) as follows.  It fixes literals, \(\bottom\), and
\(\neg\bottom\); it sends \(Qx\,\varphi\) to
\(Qx\,\operatorname{pren}(\varphi)\); and if
\[
 \operatorname{pren}(\varphi)=Q_1y_1\cdots Q_ky_k\,\varphi_0,
 \qquad
 \operatorname{pren}(\psi)=R_1z_1\cdots R_mz_m\,\psi_0,
\]
then, for \(\circ\in\{\land,\lor\}\),
\[
 \operatorname{pren}(\varphi\mathbin\circ\psi)
 =Q_1y_1\cdots Q_ky_kR_1z_1\cdots R_mz_m
   (\varphi_0\mathbin\circ\psi_0).
\]
This map selects a concrete prefix ordering; it is not the syntactic property
of already being prenex, and the distinction is needed by the displayed
clausification pipeline.
The output \(\operatorname{pren}(\varphi)\) is rectified and prenex,
\(\operatorname{pren}(\varphi)\Leftrightarrow\varphi\); its matrix is
quantifier-free NNF;
its free-variable and all-variable sets are unchanged; and its prefix lists
exactly the quantifier occurrences of \(\varphi\), with the same quantifiers
and variables and no new occurrences.
\end{hblemma}

\begin{hbtheorem}{thm:p02-fol-prenex}{proved=annex:p2.fol.nf.4}{Prenex normal form}
Every first-order formula is logically equivalent to a rectified prenex
formula in the same signature.
\end{hbtheorem}

The repaired chain is therefore
\[
 \varphi\longmapsto\operatorname{nnf}(\varphi)
 \longmapsto \operatorname{rect}\bigl(\operatorname{nnf}(\varphi),
   \operatorname{Var}(\operatorname{nnf}(\varphi))\bigr)
 \longmapsto \operatorname{pren}\bigl(
   \operatorname{rect}(\operatorname{nnf}(\varphi),\ldots)\bigr).
\]
Each arrow is a metalinguistic transformation, not an object-language
connective.  R-007 is retained in the individual lemmas rather than hidden
behind an appeal to propositional NNF.

\section{Skolemization}
\label{sec:p02-fol-skolemization}

\begin{hbdefinition}{def:p02-skolemization}{Skolemization}
\index{Skolemization}
Let \(\sigma=Q_1x_1\cdots Q_nx_n\,\chi\) be a rectified prenex
\(\mathcal L\)-sentence.  Traverse its prefix from left to right.  At each
existential occurrence \(\exists x_i\), list the variables
\(u_1,\ldots,u_m\) bound by preceding universal occurrences, introduce a
new \(m\)-ary function symbol \(f_i\), replace the occurrences bound by
\(\exists x_i\) with \(f_i(u_1,\ldots,u_m)\), and delete that existential.
Retain every universal.  Distinct existential occurrences receive distinct
new symbols; \(m=0\) produces a constant.  The resulting universal sentence
in the expanded signature is \(\operatorname{Sk}(\sigma)\).
\end{hbdefinition}

\begin{hbtheorem}{thm:p02-skolem-equisatisfiable}{proved=annex:p2.fol.nf.5}{Expansion--reduct theorem for Skolemization}
Let \(\sigma\) be a rectified prenex \(\mathcal L\)-sentence.
\begin{enumerate}
  \item Assuming the Axiom of Choice, every \(\mathcal L\)-model of
        \(\sigma\) has an expansion satisfying
        \(\operatorname{Sk}(\sigma)\).
  \item Without any choice principle, the \(\mathcal L\)-reduct of every
        model of \(\operatorname{Sk}(\sigma)\) satisfies \(\sigma\).
\end{enumerate}
Consequently, under AC, the two sentences are equisatisfiable across the
expansion--reduct relation.
\end{hbtheorem}

\begin{hbwarning}{def:p02-skolem-not-equivalence}{Skolemization boundary}
Skolemization changes the signature and names witness functions absent from
the input.  Theorem~\ref{thm:p02-skolem-equisatisfiable} states an
expansion--reduct relation and equisatisfiability, not
\(\sigma\Leftrightarrow\operatorname{Sk}(\sigma)\).  AC is used only to
choose witness functions in the expansion direction; the reduct direction
is choice-free.
\end{hbwarning}

\section{The clausal-form pipeline}
\label{sec:p02-fol-clausal-pipeline}

\begin{hbtheorem}{thm:p02-fol-clausal-form}{cited=BaazEglyLeitsch2001}{Clausal-form theorem}
For every first-order sentence \(\sigma\), there is a finite set \(S\) of
universally read clauses in a finite expansion of its signature such that,
assuming AC, \(\sigma\) is satisfiable iff \(S\) is satisfiable.  The reverse
direction from an \(S\)-model to a model of \(\sigma\) uses only reducts and
does not require choice.  The transformation is the composition in
Table~\ref{tab:p02-fol-normal-form-pipeline}; only the Skolem step weakens
same-signature logical equivalence to expansion--reduct equisatisfiability.
\end{hbtheorem}

\begin{landscape}
\begingroup\scriptsize\setlength{\tabcolsep}{3pt}
\begin{longtable}{@{}p{.30in}p{1.05in}p{1.65in}p{1.90in}p{2.00in}@{}}
\caption{First-order normal-form comparison and clausal pipeline (apparatus A-034).}
\label{tab:p02-fol-normal-form-pipeline}\\
\toprule
Step & Output & Existence and scope & Preservation claim & Worst-case size and new data \\ \midrule
\endfirsthead
\toprule Step & Output & Existence and scope & Preservation claim & Worst-case size and new data \\ \midrule
\endhead
1 & NNF & Exists for every formula in the same signature. & Formula
equivalence; exactly the same free and all-variable sets. & Linear in the
formula tree; no new symbols. \\
2 & Rectified NNF & Exists because the variable supply is countably infinite;
fresh binders are chosen outside a finite forbidden set. & Alpha-equivalence,
hence formula equivalence; same constructor tree. & Linear apart from the
encoding of fresh variable names; no new nonlogical symbols. \\
3 & Prenex NNF & Exists for every rectified NNF formula by the side-conditioned
laws in Table~\ref{tab:p02-quantifier-equivalences}. & Formula equivalence;
same free and all-variable sets and the same quantifier occurrences. & Linear
in the formula tree; only the prefix is rearranged. \\
4 & Universal Skolem sentence & Exists for every rectified prenex sentence in a
finite signature expansion; AC supplies witness functions in the forward
model-expansion direction. & Expansion--reduct equisatisfiability under AC,
not same-signature formula equivalence. & At most quadratic under literal term
printing because a new Skolem term may list all preceding universal variables;
new function symbols are introduced. \\
5 & Matrix in CNF & Exists for every quantifier-free matrix by propositional
distribution. & Formula equivalence in the expanded signature. & Direct
distribution can be exponential; Tseitin conversion is only equisatisfiable
and introduces symbols (see Part~2.1). \\
6 & Matrix in DNF & Exists for every quantifier-free matrix by dual
propositional distribution. & Formula equivalence in the expanded signature. &
Direct distribution can be exponential; DNF is a comparison form, not the
input used by resolution. \\
7 & Clause set & Exists from a universal sentence whose matrix is in CNF. & The
set and the universal closures of its clauses have exactly the same models as
that universal sentence. & Linear once the CNF matrix is present; erases the
written universal prefix and records a finite literal set per conjunct. \\
\bottomrule
\end{longtable}
\endgroup
\end{landscape}

\begin{figure}[htbp]
\centering
\begin{tikzpicture}[
  font=\small,
  node distance=8mm,
  box/.style={draw,rounded corners,align=center,inner sep=4pt,
              minimum width=1.55in},
  >=Stealth]
  \node[box] (in) {sentence\\\(\sigma\)};
  \node[box,below=of in] (nnf) {NNF};
  \node[box,below=of nnf] (rect) {rectified NNF};
  \node[box,below=of rect] (pren) {rectified prenex NNF};
  \node[box,below=of pren] (sk) {universal Skolem sentence};
  \node[box,below=of sk] (cnf) {universal sentence\\with CNF matrix};
  \node[box,below=of cnf] (clauses) {clause set \(S\)};
  \draw[->] (in) -- node[right=8mm,align=left,font=\footnotesize]
    {formula equivalence} (nnf);
  \draw[->] (nnf) -- node[right=8mm,align=left,font=\footnotesize]
    {alpha-equivalence} (rect);
  \draw[->] (rect) -- node[right=8mm,align=left,font=\footnotesize]
    {formula equivalence} (pren);
  \draw[->] (pren) -- node[right=8mm,align=left,font=\footnotesize]
    {equisatisfiable\\expanded signature} (sk);
  \draw[->] (sk) -- node[right=8mm,align=left,font=\footnotesize]
    {matrix CNF\\formula equivalence} (cnf);
  \draw[->] (cnf) -- node[right=8mm,align=left,font=\footnotesize]
    {universal-closure convention} (clauses);
\end{tikzpicture}
\caption{Normal-form pipeline (apparatus A-118).  The Skolem arrow records
expansion--reduct equisatisfiability rather than same-signature formula
equivalence; the last arrow relates a sentence to a clause set through the
universal closures of its members, not through formula-string equivalence.}
\label{fig:p02-fol-normal-form-pipeline}
\end{figure}

This theorem is the S-level composition required by D-P3-012.  Resolution
objects and their universal-closure reading are defined in the next chapter;
the table does not identify a finite clause set with a single formula string.
For the quantifier equivalences and the clause-conversion pipeline,
respectively, see \cite{Enderton2001,BaazEglyLeitsch2001}.
